\bookStart{Against Wen}[Wið Wennan]
\setBookCode{Wen}

\begin{flushright}%
\textbf{Dating:} ca.~1150 \ce

\textbf{Meter:} \Fornyrdislag%
\end{flushright}

\section{Introduction}

TODO: Introduction.

British Library MS Royal 4.A.xiv, fol.~106v.

An exorcistic galder.  The text has early Middle English features like the turning of old \emph{f} to \emph{v}.

\sectionline

\section{Text}

\bvg\bva%
\alst{W}enne, \alst{W}enne \hld\ \alst{W}enchi-chenne, &
hér ne scealt þú \alst{t}imbrien, \hld\ ne nenne \alst{t}ún habben, &
ac þú scealt \alst{n}ort heonene \hld\ tó þan \alst{n}ihgan berhge, &
þer þú havest \alst{e}rming, \hld\ \alst{é}nne bróþer. &
Hé þe sceal \alst{l}egge \hld\ \alst{l}eaf et héafde. &
Under \alst{f}ót wolmes, \hld\ under \alst{v}eþer earnes, &
under \alst{ea}rnes cléa, \hld\ \alst{á} þú ge·weornie. &\eva

\bvb TODO: translation\evb\evg


\bvg\bva[8]%
\alst{C}linge þú al-swa \hld\ \alst{c}ol on heorþe, &
\alst{s}cring þú al-swa \hld\ \alst{s}cerne awage, &
and \alst{w}eorne \hld\ al-swa \alst{w}eter on anbre.
Swa \alst{l}itel þú ge·wurþe \hld\ alswa \alst{l}inset-corn, &
and miccli lesse al-swa ánes hand-wurmes hupeban, &
and al-swa litel þú ge·wurþe þet þú ná-wiht ge·wurþe.\eva

\bvb TODO: translation\evb\evg

\sectionline
